% page 113 du fac-simile (image p-112 du scan)
% bloc 172.7 x 237.7 mm ; marge gauche 29.3 mm
\pgc{7.620mm}{0.000mm}{
\l{\cel{57}lO3 bR}
\l{}
\l{}
\l{\sou{dermo}. (\sou{anat.}) Tisuo-fibroza ye qua konsistas la strato maxim dika}
\l{\cel{3}e maxim profunde-situita di pelo, e quan kovras membrano surfaca-}
\l{\cel{3}la (epidermo). - DEFIS.}
\l{}
\l{\sou{dermatogeno}. (\sou{anat.}) Tisuo yuna qua, ye la somito di la stipo o di}
\l{\cel{3}la radiko, produktas la epidermo o la \textquotedbl{}kofio\textquotedbl{}.}
\l{}
\l{\sou{deroga}r (\sou{netrans.}\cel{1}de). (\sou{Yurocienco).}\cel{2}Eceptar, violacar provizore,}
\l{\cel{3}ula lego o regulo. - DEFIS.}
\l{}
\l{\sou{dervisho}. (\sou{Che la Mohamedisti}).\cel{2}Quaza monako. - DEFIRS.}
\l{}
\l{\sou{\textquotedbl{}des\textquotedbl{}}.\cel{2}(\sou{Muziko}). La noto duesma di la \textquotedbl{}do\textquotedbl{}-gamo bemoloza.}
\l{}
\l{\sou{des-}\cel{4}La antonimo di (to quon indikas la radiko).}
\l{}
\l{\sou{desegnar}.\cel{2}(\sou{trans, per, sur, e}\cel{1}c.)I.Riproduktar, per krayono o}
\l{\cel{3}plumo, la formo di la objekti. II. Igar videbla la konturi.}
\l{\cel{3}- EFIS.}
\l{}
\l{\sou{desero}. La kozi servata laste, lor repasto : fromajo, frukti,}
\l{\cel{3}kuko, e c. - DEFR.}
\l{}
\l{\sou{desertar}. (\sou{trans, e netrans}).\cel{2}Abandonar loko (precipue aludante}
\l{\cel{3}armeano) o partiso, e c. quin onu devas ne livar. - DEFIRS.}
\l{}
\l{\sou{deservar}. (\sou{trans}.) (en\cel{1}\sou{templo, kirko, sinagogo}). Agar la servado}
\l{\cel{3}religiala. - DEFS.}
\l{}
\l{\sou{desikatoro}.I.(\sou{kemio}). Aparato sikigera.II. (\sou{tekn}.)Aparato en qua}
\l{\cel{3}onu pozas diversa texo-materii por determinar la pezo vera di la}
\l{\cel{3}vari en paki. - DEFIRS.}
\l{}
\l{\sou{deskriptar}. (\sou{trans.}) Reprezentar skribe o parole (persono, kozo)}
\l{\cel{3}pri lua traiti extera.-- EFIS.}
\l{}
\l{\sou{despenso}. Loko ube onu inkluzas la provizuri (en kolegio, farmo-}
\l{\cel{3}domo, e c.) - FIS.}
\l{}
\l{\sou{despitar}. (\sou{netrans., pri, pro,}\cel{1}de). Iracar pro desprizeso, des-}
\l{\cel{3}egardeso, o pro ke altru preferesas a ni, e chagrenas pri lo.}
\l{\cel{3}- EFIS.}
\l{}
\l{\sou{despoto}. Monarko absolutista qua guvernas segun sua arbitrio.}
\l{\cel{3}- DEFIRS.}
\l{}
\l{\sou{destinar}. (\sou{trans., ad, por)}.\cel{2}I. Fixigar anticipe (ad ulu) kom}
\l{\cel{3}esonta lua loto. - II. Fixigar anticipe kom uzota por ta o ca}
\l{\cel{3}skopo. - DEFIS.}
\l{}
\l{\sou{destruktar}. (\sou{trans}.) I. Desfacar (lo facita). - II. Ruinar tote.}
\l{\cel{3}- DEFIRS.}
}
